\subsection{自由族，基}

设 A 为一个环，T 为一个集合，并考虑 A-模 \(\mathbf{A}_{s}^{\mathrm{(T)}}\) 。根据定义，它是一族 A-模 \((\mathbf{M}_{t})_{t\in\mathbb{T}}\) （这些 A-模都等于 \(A_{s}\) ）的外直和，且对所有 \(t\in T\) 存在典范单射 \(j_{t}:A_{s}\to A_{s}^{\mathrm{(T)}}\) （第 6 号）。我们记 \(j_{t}(1) = e_{t}\) ，使得 \(e_{t} = (\delta_{tt'})_{t' \in \mathbb{T}}\) ，其中 \(\delta_{tt'}\) 当 \(t' \neq t\) 时等于 0，当 \(t' = t\) 时等于 1（“克罗内克符号”； \((t, t') \mapsto \delta_{tt'}\) 恰为 \(T \times T\) 的对角线的特征函数）；于是每个 \(x = (\xi_{t})_{t \in T} \in A_{s}^{(T)}\) 可唯一地写成： \(x = \sum_{t \in T} \xi_{t} e_{t}\) 。映射 \(\phi: t \mapsto e_{t}\) （ \(T\) 到 \(A_{s}^{(T)}\) 中的映射）称为典范映射；若 \(A\) 非零，则它是单射。我们将看到，有序对 \((A_{s}^{(T)}, \phi)\) 是某个泛映射问题的一个解（集合论，IV，§3，第 1 号）。

命题 17. 对每个 A-模 E 和每个映射 \(f: \mathbf{T} \to \mathbf{E}\) ，存在且仅存在一个 A-线性映射 \(g: \mathbf{A}_{s}^{(\mathrm{T})} \to \mathbf{E}\) ，使得 \(f = g \circ \phi\) 。

条件 \(f = g \circ \phi\) 表示 \(g(e_t) = f(t)\) 对所有 \(t \in \mathbf{T}\) 成立，这等价于 \(g(\xi e_t) = \xi f(t)\) 对所有 \(\xi \in \mathbf{A}\) 和所有 \(t \in \mathbf{T}\) 成立，并且也表示 \(g \circ j_t\) 是线性映射 \(\xi \mapsto \xi f(t)\) （ \(\mathbf{A}_s\) 到 \(\mathbf{E}\) 中的映射）对所有 \(t \in \mathbf{T}\) 成立；因此该命题是第 6 号命题 6 的一个特例。

线性映射 g 称为由 E 中元素的族 \((f(t))_{t\in\mathbb{T}}\) 所确定；根据定义

\[
g \left(\sum_ {t \in \mathrm{T}} \xi_ {t} e _ {t}\right) = \sum_ {t \in \mathrm{T}} \xi_ {t} f (t).\tag{43}
\]

核 \(\mathbf{R}\) （即 \(g\) 的核）是 \((\xi_t) \in \mathbf{A}_s^{(\mathrm{T})}\) （满足 \(\sum_{t} \xi_t f(t) = 0\) 者）组成的集合；有时称模 \(\mathbf{R}\) 为族 \((f(t))_{t \in \mathbb{T}}\) 的诸元素之间线性关系构成的模。正合列

\[
0 \longrightarrow \mathrm{R} \longrightarrow \mathrm{A} _ {s} ^ {(\mathrm{T})} \stackrel {\mathbf {g}} {\longrightarrow} \mathrm{E}\tag{44}
\]

称为由族 \((f(t))_{t\in\mathbb{T}}\) 所确定。

推论 1. 设 T、T' 为两个集合，g: T → T' 为一个映射。则存在且仅存在一个 A-线性映射 f: A \(^{(T)}\) → A \(^{(T')}\) ，使得下图交换

\[
\begin{array}{c c c} \mathrm{T} & \stackrel {{g}} {{\longrightarrow}} & \mathrm{T} ^ {\prime} \\ \phi \Big \downarrow & & \Big \downarrow \phi^ {\prime} \\ \mathrm{A} ^ {(\mathrm{T})} & \stackrel {{f}} {{\longrightarrow}} & \mathrm{A} ^ {(\mathrm{T} ^ {\prime})} \end{array}
\]

其中 \(\phi\) 和 \(\phi'\) 是典范映射。

只需将命题 17 应用于复合映射 \(\mathbf{T} \xrightarrow{\pmb{g}} \mathbf{T}' \xrightarrow{\phi'} \mathbf{A}^{(\mathbf{T}')}\) 。

推论 2. 要族 \((a_{t})_{t\in \mathbb{T}}\) （ A-模 \(\mathbf{E}\) 中元素的一个族）是 \(\mathbf{E}\) 的一个生成系，必须且只需该族所确定的线性映射 \(\mathbf{A}_{s}^{(T)}\to \mathbf{E}\) 是满射。

这只是第 7 号命题 9 的另一种表述方式。

定义 10. 族 \((a_{t})_{t\in \mathbf{T}}\) （ A-模 \(\mathbf{E}\) 中元素的一个族）称为自由族（相应地，称为 \(\mathbf{E}\) 的一组基），如果该族所确定的线性映射 \(\mathbf{A}_{s}^{(\mathbf{T})}\to \mathbf{E}\) 是单射（相应地，是双射）。若一个模具有一组基，则称它为自由模。

特别地，如果交换群 G （以加法记号书写）是一个自由 Z-模，就称它是自由的（参见 I，§7，第 7 号）。

定义 10 连同命题 17 的推论 2 表明，A-模 E 的一组基是 E 的一个自由生成族。因此，E 中元素的一个自由族就是它所生成的子模的一组基。

根据定义，A-模 \(\mathbf{A}_s^{(\mathrm{T})}\) 是自由的，且族 \((e_t)_{t\in \mathbb{T}}\) 是该 A-模的一组基（称为典范基）。当 \(\mathbf{A}\neq \{0\}\) 时， \(\mathbf{T}\) 常与 \(e_t\) 的集合（经由典范双射 \(t\mapsto e_t\) ）恒同；这相当于把 \(\sum_{t\in \mathbb{T}}\xi_t.t\)

（而不是 \(\sum_{t\in T}\xi_{t}a_{t}\) ）写出来表示 \(A_{s}^{(T)}\) 的元素。当采用这一约定时， \(A_{s}^{(T)}\) 的元素称为 T 的元素的（系数在 A 中的）形式线性组合。

定义 10 和命题 17 立即给出以下结果：

推论 3. 设 E 是一个自由 A-模， \((a_{t})_{t\in \mathbb{T}}\) 是 E 的一组基，F 是一个 A-模，且 \((b_{t})_{t\in \mathbb{T}}\) 是 F 中元素的一个族。则存在且仅存在一个线性映射 \(f:\mathbf{E}\to \mathbf{F}\) ，使得

\[
f (a _ {t}) = b _ {t} \text {   for   all   } t \in \mathrm{T}.\tag{45}
\]

要 \(f\) 是单射（相应地，满射），必须且只需 \((b_t)\) 是 \(\mathbf{F}\) 中的一个自由族（相应地，是 \(\mathbf{F}\) 的一个生成系）。

当族 \((a_{t})_{t\in \mathbb{T}}\) 不是自由族时，称它为相关的。定义 10 也可以表述如下：称族 \((a_{t})_{t\in \mathbb{T}}\) 是自由的，意思是关系式 \(\sum_{t\in \mathbb{T}}\lambda_t a_t = 0\) （其中族 \((\lambda_t)\) 的支撑有限）蕴含 \(\lambda_{t} = 0\) 对所有 \(t\in \mathbb{T}\) 成立；称 \((a_{t})_{t\in \mathbb{T}}\) 是 E 的一组基，意思是每个 \(x\in \mathbb{E}\) 都能以一种且仅一种方式写成 \(x = \sum_{t\in \mathbb{T}}\xi_t a_t\) 的形式；这时，对所有 \(t\in \mathbb{T}\) ，称 \(\xi_{t}\) 为 \(x\) 的下标为 \(t\) 的分量（或坐标）（在基 \((a_{t})\) 下）；映射 \(x\to \xi_t\) （从 E 到 \(\mathbf{A}_s\) 的映射）是线性的。

假设 \(\mathbf{A} \neq \{0\}\) ；于是在 A-模 \(\mathbf{E}\) 中，自由族 \((a_t)_{t \in \mathbb{T}}\) 的指标不同的两个元素本身也不相同：因为若 \(a_{t'} = a_{t''}\) （其中 \(t' \neq t''\) ），则 \(\sum_{t \in \mathbb{T}} \lambda_t a_t = 0\) ，其中 \(\lambda_{t'} = 1\) ， \(\lambda_{t''} = -1\) ，且 \(\lambda_t = 0\) 对 \(\mathbb{T}\) 中异于 \(t'\) 与 \(t''\) 的元素成立。子集 \(\mathbb{S}\) （ \(\mathbb{E}\) 的一个子集）称为自由子集（相应地， \(\mathbb{E}\) 的一个基），如果由 \(\mathbb{S}\) 到其自身的恒同映射所定义的族是自由的（相应地，是 \(\mathbb{E}\) 的一个基）；这时，由指标集到 \(\mathbb{S}\) 上的一个双射映射所定义的每个族都是自由的（相应地，是一个基）。 \(\mathbb{E}\) 的一个自由子集的元素也称为线性无关的。

如果 E 的一个子集不是自由子集，则称它为相关的，或称其为一个相关系，它的元素称为线性相关的。

自由子集的每个子集都是自由的；特别地，空子集是自由的，并且是 E 的子模 \(\{0\}\) 的一个基。

命题 18. 族 \((a_{t})_{t\in \mathbb{T}}\) （模 \(\mathbf{E}\) 的一个族）是自由的，必须且只需 \((a_{t})_{t\in \mathbb{T}}\) 的每个有限子族都是自由的。

这直接从定义得出。

命题 18 表明，E 的自由子集的集合按包含关系排序是归纳的（集合论，III，§2，第 4 号）；由于它非空（因为 \(\varnothing\) 属于它），根据佐恩引理它有一个极大元素 \((a_{i})_{i\in I}\) 。由此可得（若 \(\mathbf{A}\neq \{0\}\) ）：对所有 \(x\in \mathbf{E}\) ，存在 A 的一个元素 \(\mu \neq 0\) 以及 A 的元素的一个族 \((\xi_{i})\) ，使得 \(\mu x = \sum_{i}\xi_{i}a_{i}\) （参见 §7，第 1 号）。

命题 19. 设 E 是一个 A-模，且是子模的一个族 \((\mathbf{M}_{\lambda})_{\lambda \in L}\) 的直和。如果对每个 \(\lambda \in L\) ， \(S_{\lambda}\) 是 \(M_{\lambda}\) 的一个自由子集（相应地，生成集、基），则 \(S = \bigcup_{\lambda \in L} S_{\lambda}\) 是 E 的一个自由子集（相应地，生成集、基）。

该命题由定义以及关系式 \(\mathbf{A}_{s}^{(\mathbf{S})} = \bigoplus_{\lambda \in \mathbf{L}} \mathbf{A}_{s}^{(\mathbf{S}_{\lambda})}\) （直和的结合律，参见第 6 号）得出。

注记。(1) 根据定义 10，若 \(\mathbf{A} \neq \{0\}\) 且 \((a_{t})_{t \in I}\) 是一个自由族，则任何元素 \(a_{k}\) 都不可能等于那些 \(a_{t}\) （其指标 \(t \neq k\) ）的线性组合。但反过来，满足此条件的一个族 \((a_{t})\) 未必是自由族。例如，设 \(\mathbf{A}\) 是一个整环， \(a, b\) 是两个不同的非零元素；在 \(\mathbf{A}\) （视为 A-模）中， \(a\) 和 \(b\) 构成一个相关系，因为 \((-b)a + ab = 0\) 。但一般而言，不存在元素 \(x \in \mathbf{A}\) 使得 \(a = xb\) 或 \(b = xa\) （然而可参见 §7，第 1 号的注记）。

模 E 的一个元素 x 称为自由的，如果 \(\{x\}\) 是一个自由子集，也就是说，如果关系式 \(\alpha x = 0\) 蕴含 \(\alpha = 0\) 。自由子集的每个元素都是自由的，特别地，当 \(A \neq \{0\}\) 时，0 不可能属于任何自由子集。

注记。(2) 一个自由模可以含有 \(\neq0\) 的非自由元素：例如，A-模 \(A_{s}\) 是自由的，但 A 中的右零因子不是 \(A_{s}\) 的自由元素。

\begin{enumerate}
\def\labelenumi{(\arabic{enumi})}
\setcounter{enumi}{2}
\item
  在加法群 \(\mathbf{Z} / (n)\) （ \(n\) 为一个整数 \(\geqslant 2\) ，且被视为一个 \(\mathbf{Z}\) -模）中，没有元素是自由的，从而更不用说 \(\mathbf{Z} / (n)\) 不是自由模。
\item
  可能发生这样的情形：一个 A-模的每个元素 \(\neq0\) 都是自由的，而该模本身却不是自由的。例如，有理数域 Q 作为 Z-模就具有这一性质，因为 Q 的两个元素 \(\neq0\) 总是构成一个相关系，因此 \(\mathbf{Q}\) 的基只能含一个元素 \(a\) ；但 \(\mathbf{Q}\) 的元素并非都具有 \(na\) （ \(n \in \mathbf{Z}\) ）的形式（参见 VII，§3）。
\end{enumerate}

命题 20. 每个 A-模 E 都同构于一个自由 A-模的商模。

如果 T 是 E 的一个生成集，则存在一个满射线性映射 \(\mathbf{A}_{s}^{(\mathrm{T})}\to\mathbf{E}\) （命题 17 的推论 2），且若 R 是这个映射的核，则 E 同构于 \(\mathbf{A}_{s}^{(\mathrm{T})}/\mathbf{R}\) 。

特别地，我们可以取 \(\mathbf{T} = \mathbf{E}\) ；这时存在一个满射线性映射 \(\mathbf{A}_s^{(\mathbf{E})} \to \mathbf{E}\) ，称为典范映射。

特别地，说一个 A-模 E 是有限生成的（第 7 号），意味着它同构于一个具有有限基的自由 A-模的商模，或者等价地，存在如下形式的正合列

\[
\mathrm{A} _ {s} ^ {n} \rightarrow \mathrm{E} \rightarrow 0 (n \text {   an   integer   } > 0).
\]

注意，如果 \(A \neq \{0\}\) ，有限生成的自由模 E 的每个基都必然是有限的，因为若 S 是一个有限生成系而 B 是 E 的一个基，则 S 的每个元素都是 B 中有限个元素的线性组合，且若 \(B'\) 是 B 中所有以这种方式出现在 S 的元素的表达式中的元素组成的集合，则 \(B'\) 是有限的，而每个 \(x \in E\) 都是 \(B'\) 中元素的线性组合，因此 \(B' = B\) 。

命题 21. A-模的每个正合列

\[
0 \longrightarrow \mathrm{G} \stackrel {g} {\longrightarrow} \mathrm{E} \stackrel {f} {\longrightarrow} \mathrm{F} \longrightarrow 0
\]

其中 \(\mathbf{F}\) 是一个自由 A-模，是分裂的（第 9 号）。确切地说，如果 \((b_{\lambda})_{\lambda \in \mathbf{L}}\) 是 \(\mathbf{F}\) 的一组基，且对每个 \(\lambda \in \mathbf{L}\) ， \(a_{\lambda}\) 是 \(\mathbf{E}\) 的一个元素，使得 \(f(a_{\lambda}) = b_{\lambda}\) ，则族 \((a_{\lambda})_{\lambda \in \mathbf{L}}\) 是自由的，并生成与 \(g(\mathbf{G})\) 互补的子模。

存在且仅存在一个线性映射 \(h: \mathbf{F} \to \mathbf{E}\) ，使得 \(h(b_{\lambda}) = a_{\lambda}\) 对所有 \(\lambda \in \mathbf{L}\) 成立（命题 17 的推论 3）。由于 \(h\) 是与 \(f\) 相伴的线性截面，该命题由 I，§4，第 9 号的命题 15 得出。

注记。(5) 设 \((a_{i})_{1\leqslant i\leqslant n}\) 是 A-模 E 的一组基，且设 \((b_{i})_{1\leqslant i\leqslant n}\) 是 E 中由下列关系式给出的一个元素族：

\[
b _ {i} = \lambda_ {1 i} a _ {1} + \dots + \lambda_ {i i} a _ {i} (1 \leqslant i \leqslant n)\tag{46}
\]

其中 \(\lambda_{ii}\) 在 A 中可逆；则 \((b_i)_{1\leqslant i\leqslant n}\) 是 E 的一组基。只需对 \(n\) 作归纳论证，因为该命题对 \(n = 1\) 是显然的。若 \(\mathbf{E}'\) 是 E 的由族 \((a_i)_{1\leqslant i\leqslant n - 1}\) 生成的子模，则由归纳假设可得 \((b_i)_{1\leqslant i\leqslant n - 1}\) 是 \(\mathbf{E}'\) 的一组基；另一方面，由 (46) 可得，如果 \(\mu b_n\in \mathbf{E}'\) （其中 \(\mu \in \mathbf{A}\) ），那么也有 \(\mu \lambda_{nn}a_n\in \mathbf{E}'\) ，从而 \(\mu = 0\) ，因为 \(\lambda_{nn}\) 是可逆的。因此族 \((b_i)_{1\leqslant i\leqslant n}\) 是自由的，并且由于

\[
a _ {n} = - \lambda_ {n n} ^ {- 1} \lambda_ {1 n} a _ {1} - \dots - \lambda_ {n n} ^ {- 1} \lambda_ {n - 1, n} a _ {n - 1} + \lambda_ {n n} ^ {- 1} b _ {n}
\]

由此可以看出 \((b_{i})_{1\leqslant i\leqslant n}\) 是 E 的一个生成系，这就完成了证明。这一结果容易推广到指标集 I 为良序集的族 \((a_{i})_{i\in I}\) 。
% semantic-fragments: legacy-input-seam
\relax{}
